\par\Needspace{6\baselineskip}\medskip
\hypertarget{native-211-article156}{}\pdfbookmark[2]{156.}{article-156}
\begin{center}156.\end{center}
Unser \textit{zweites} Beispiel nehmen wir von Pallas. Die folgenden in Mailand angestellten Beobachtungen entnehmen wir der gelehrten Korrespondenz des berühmten \textsc{de Zach}, Bd.\ XIV. S.\ 90.
\begin{center}
\begingroup\footnotesize\setlength{\tabcolsep}{4pt}\renewcommand{\arraystretch}{1.2}
\begin{tabular}{|l|r|l|}\hline
Mittlere Mailänder Zeit & Scheinb.\ Rektasz. & Scheinb.\ Dekl.\\\hline
1805 Nov.\ 5.\ $14^{\mathrm h}\,14^{\mathrm m}\,4^{\mathrm s}$ & $78^{0}\,20'\,37''{,}8$ & $27^{0}\,16'\,56''{,}7$ südl.\\
\hspace{2.4em}Dez.\ 6.\ $11\quad51\quad27$ & $73\quad8\quad48{,}8$ & $32\quad52\quad44{,}3$\\
1806 Jan.\ 15.\ $8\quad50\quad36$ & $67\quad14\quad11{,}1$ & $28\quad38\quad8{,}1$\\\hline
\end{tabular}\endgroup
\end{center}

\Needspace{9\baselineskip}
Hier nehmen wir anstelle der Ekliptik den Äquator als Grundebene und führen die Rechnung so aus, als wäre die Bahn noch völlig unbekannt. Zunächst entnehmen wir den Sonnentafeln für die gegebenen Zeiten Folgendes:
\begin{center}
\begingroup\footnotesize\setlength{\tabcolsep}{4pt}\renewcommand{\arraystretch}{1.2}
\begin{tabular}{|l|r|r|r|}\hline
& \raisebox{0pt}[\dimexpr\height+2pt\relax]{\shortstack{Länge der Sonne vom\\mittl.\ Äquinoktium}} & Abstand von der Erde & \shortstack{Breite der\\Sonne}\\\hline
Nov.\ 5 & $223^{0}\,14'\,7''{,}61$ & $0{,}990\,4311$ & $+0''{,}59$\\
Dez.\ 6 & $254\quad28\quad42{,}59$ & $0{,}984\,6753$ & $+0{,}12$\\
Jan.\ 15 & $295\quad5\quad47{,}62$ & $0{,}983\,8153$ & $-0{,}19$\\\hline
\end{tabular}\endgroup
\end{center}

Durch Hinzufügen der Präzessionen $+7''{,}59$, $+3''{,}36$, $-2''{,}11$ reduzieren wir die Längen der Sonne auf den Beginn des Jahres 1806. Daraus leiten wir mit der mittleren Schiefe $23^{0}\,27'\,53''{,}53$ und unter gebührender Berücksichtigung der Breiten die Rektaszensionen und Deklinationen ab. So erhalten wir
\SourceLeaf{212}
\begin{center}
\begingroup\small\setlength{\tabcolsep}{5pt}\renewcommand{\arraystretch}{1.2}
\begin{tabular}{|l|r|l|}\hline
& Rektaszension der Sonne & Deklination der Sonne\\\hline
Nov.\ 5 & $220^{0}\,46'\,44''{,}65$ & $15^{0}\,49'\,43''{,}94$ südl.\\
Dez.\ 6 & $253\quad9\quad23{,}26$ & $22\quad33\quad39{,}45$\\
Jan.\ 15 & $297\quad2\quad51{,}11$ & $21\quad8\quad12{,}98$\\\hline
\end{tabular}\endgroup
\end{center}

Diese Positionen beziehen sich auf den Erdmittelpunkt und sind daher durch Hinzufügen der Parallaxe auf den Beobachtungsort zu reduzieren, da man die Positionen des Planeten nicht von der Parallaxe befreien kann. Die bei dieser Rechnung zu verwendenden Rektaszensionen des Zenits stimmen mit den Rektaszensionen des Planeten überein (da die Beobachtungen im Meridian selbst angestellt wurden). Die Deklination ist in jedem Fall die Polhöhe $=45^{0}\,28'$. Daraus ergeben sich die folgenden Zahlen:
\begin{center}
\begingroup\footnotesize\setlength{\tabcolsep}{3pt}\renewcommand{\arraystretch}{1.2}
\begin{tabular}{|l|r|l|r|}\hline
& Rektasz.\ der Erde & Deklination der Erde & Log.\ Abst.\ von der Sonne\\\hline
Nov.\ 5 & $40^{0}\,46'\,48''{,}51$ & $15^{0}\,49'\,48''{,}59$ nördl. & $9{,}995\,8375$\\
Dez.\ 6 & $73\quad9\quad23{,}26$ & $22\quad33\quad42{,}83$ & $9{,}993\,3099$\\
Jan.\ 15 & $117\quad2\quad46{,}09$ & $21\quad8\quad17{,}29$ & $9{,}992\,9259$\\\hline
\end{tabular}\endgroup
\end{center}

Die beobachteten Orte von Pallas sind von der Nutation und der Aberration der Fixsterne zu befreien und dann durch Hinzufügen der Präzession auf den Beginn des Jahres 1806 zu reduzieren. Unter diesen Bezeichnungen sind die folgenden Korrekturen auf die beobachteten Positionen anzuwenden:
\begin{center}
\begingroup\footnotesize\setlength{\tabcolsep}{4pt}\renewcommand{\arraystretch}{1.2}
\begin{tabular}{|l|r|r||r|r||r|r|}\hline
& \multicolumn{2}{c||}{Beobachtung I} & \multicolumn{2}{c||}{Beobachtung II} & \multicolumn{2}{c|}{Beobachtung III}\\\hline
& Rektasz. & Dekl. & Rektasz. & Dekl. & Rektasz. & Dekl.\\\hline
Nutation & $-12''{,}86$ & $-3''{,}08$ & $-12''{,}21$ & $-3''{,}42$ & $-13''{,}06$ & $-3''{,}75$\\
Aberration & $-18{,}13$ & $-9{,}89$ & $-22{,}68$ & $-1{,}63$ & $-15{,}60$ & $+9{,}76$\\
Präzession & $+5{,}43$ & $+0{,}62$ & $+2{,}25$ & $+0{,}39$ & $-1{,}51$ & $-0{,}33$\\\hline
Summe & $-25{,}56$ & $-12{,}35$ & $-32{,}64$ & $-4{,}66$ & $-30{,}17$ & $+5{,}68$\\\hline
\end{tabular}\endgroup
\end{center}

So ergeben sich die folgenden Positionen von Pallas, die der Rechnung zugrunde zu legen sind:
\begin{center}
\begingroup\small\setlength{\tabcolsep}{6pt}\renewcommand{\arraystretch}{1.2}
\begin{tabular}{|r|r|r|}\hline
Mittl.\ Pariser Zeit & Rektasz. & Deklination\\\hline
Nov.\ $5{,}574\,074$ & $78^{0}\,20'\,12''{,}24$ & $-27^{0}\,17'\,9''{,}05$\\
$36{,}475\,035$ & $73\quad8\quad16{,}16$ & $-32\quad52\quad48{,}96$\\
$76{,}349\,444$ & $67\quad13\quad40{,}93$ & $-28\quad38\quad2{,}42$\\\hline
\end{tabular}\endgroup
\end{center}
